% =============================================================================
%  CHAPTER 14 -- PHASE 1: DATA ACQUISITION AND UNIVERSE CONSTRUCTION
% =============================================================================
\section{Phase 1 --- Data acquisition and universe construction}
\label{ch:phase1}

\textbf{Reproduce it:} \code{python3 scripts/data/clean.py \&\& python3 scripts/data/qa.py
\&\& python3 scripts/analysis/make\_data\_summary.py \&\& python3
scripts/plotting/render\_report\_data.py}\\
\textbf{Math used:} Section~\ref{sec:prob:returns} (returns, adjustment arithmetic).\\
\textbf{Artifacts:} \code{data/processed/universe.csv}, \code{symbols.csv},
\code{\_meta.json}; \code{results/tables/\{data\_summary,universe\_by\_year,sector\_counts\}.csv};
\code{results/figures/\{universe\_by\_year,sector\_comp\}.png};
\code{report/tables/data\_section.tex}.

\subsection{Goal}

Obtain a real, documented, corporate-action-adjusted equity panel: $\ge200$ names,
$\ge10$ years of daily data, no lookahead in construction, survivorship bias quantified
rather than ignored.

\subsection{What was done}
\label{sec:phase1:universe}

\subsubsection{Acquisition}

Two sources, one primary:

\begin{description}[leftmargin=0pt,style=nextline]
  \item[Primary --- Yahoo Finance via \code{yfinance}] daily OHLCV plus \code{Adj Close}
    for the S\&P-500-style symbol list in \code{scripts/data/sp500\_tickers.txt}
    (505 names), start 2010-01-01, retrieved 2026-09-04 on an internet-enabled machine
    and committed as one CSV per ticker under \code{data/raw/yfinance/}. \textbf{195} of
    the 505 returned a usable history, giving the effective universe.
  \item[Secondary --- NYSE/Kaggle panel via GitHub tarball] 501 U.S.\ securities,
    split-adjusted daily OHLCV, GICS sectors, 2010-01-04 to 2016-12-30, fetched in-sandbox
    by \code{download.py --source github} from
    \code{api.github.com/repos/\{owner\}/\{repo\}/tarball/\{ref\}}. Used only for
    cross-sectional breadth checks, because it is split-adjusted \emph{only} --- returns
    computed from it are price returns, missing $\sim$1.5--2.5\%/year of dividends
    (Section~\ref{sec:prob:returns}).
\end{description}

\code{download.py} writes a provenance manifest (source, URL/ref, retrieval timestamp,
SHA-256 per file) to \code{data/manifest/}, so the data is auditable and the checksums
detect silent changes. Both paths are idempotent.

\subsubsection{Cleaning}

\code{scripts/data/clean.py} ingests whichever raw source is present (preferring
\code{yfinance}) and produces a uniform long panel
\code{data/processed/universe.csv} with columns
\[
  \texttt{date},\ \texttt{ticker},\ \texttt{open},\ \texttt{high},\ \texttt{low},\
  \texttt{close},\ \texttt{adj\_close},\ \texttt{volume}.
\]
Operations, in order: drop rows with non-positive prices or non-positive volume
($0.07\%$ of the source); drop duplicate (date, ticker) pairs, last value wins; sort by
ticker then date; carry Yahoo's \code{Adj Close} through verbatim as \code{adj\_close}
(it is adjusted for \emph{both} splits and cash dividends, so returns on it are total
returns, equation~\eqref{eq:totalreturn}); write \code{symbols.csv} with name, GICS
sector, sub-industry and date-first-added from the bundled
\code{sp500\_symbols.csv}/\code{sector\_map.json}; write \code{\_meta.json} recording the
adjustment convention.

\subsubsection{Quality gate}

\code{scripts/data/qa.py} runs ten mechanical checks, all passing:

\begin{center}
\renewcommand{\arraystretch}{1.15}
\small
\begin{tabular}{lp{0.62\linewidth}}
\toprule
Result & Check \\
\midrule
PASS & non-empty --- 764{,}896 rows \\
PASS & positive \code{close} --- 0 non-positive \\
PASS & positive \code{adj\_close} --- 0 non-positive \\
PASS & no NaN \code{close} \\
PASS & no NaN \code{adj\_close} \\
PASS & no duplicate (date, ticker) --- 0 duplicates \\
PASS & no zero/negative volume \\
PASS & $\ge150$ tickers --- \textbf{195} (specification target 200) \\
PASS & $\ge250$ trading days per symbol (median) --- median \textbf{4{,}190} days \\
PASS & no excessive calendar gaps --- \textbf{38} gaps $>5$ calendar days across all symbols \\
\bottomrule
\end{tabular}
\end{center}

Per-symbol gap statistics are written to \code{results/tables/qa\_gaps.csv}
(minimum gap 1 day, maximum 5 days for essentially every name --- i.e.\ the panel is a
clean trading calendar with no missing-price holes).

\subsubsection{Engine input export}

\code{scripts/data/export\_engine\_input.py --target T --basket B1 B2 \dots} pivots the
long panel to a \emph{strictly balanced} wide matrix (any date on which any name is
missing is dropped), keeps only positive prices, and writes
\code{data/processed/engine/<T>.csv} plus a JSON sidecar recording the target, basket,
per-ticker sectors, observation count and date range. Four baskets are exported, and they
are the four used in every later phase:

\begin{center}
\renewcommand{\arraystretch}{1.2}
\begin{tabular}{llp{0.30\linewidth}rr}
\toprule
Target & Basket & Economic rationale & Bars & Span \\
\midrule
\code{ADBE} & \code{CRM}, \code{ADSK}, \code{INTU} & application software (same GICS sub-industry group) & 4{,}190 & 2010-01-04 to 2026-09-01 \\
\code{AVGO} & \code{ADI}, \code{INTC}, \code{AAPL} & semiconductors / large-cap tech hardware & 4{,}191 & same \\
\code{GS} & \code{BAC} & money-centre / large-cap banks & 4{,}190 & same \\
\code{CAT} & \code{HON} & capital goods / diversified industrials & 4{,}190 & same \\
\bottomrule
\end{tabular}
\end{center}

The drop-any-missing-date rule is a lookahead defence, not just a convenience:
forward-filling a stale price would let the strategy trade on a value that no longer
exists in the market (Section~\ref{sec:bt:lookahead}).

\subsection{Results}

\input{tables/data_section.tex}

The universe by year (Table~\ref{tab:data_summary} and
Figure~\ref{fig:universe_year}) grows from 176 names with data in 2010 to 191 in 2026 ---
not because names were added, but because later IPOs (and later index additions) enter the
panel mid-sample. The sector distribution (Figure~\ref{fig:sector}) is that of the large-cap
U.S.\ market with a 2016-vintage tilt: Consumer Discretionary 85, Industrials 69,
Information Technology 68, Financials 64, Health Care 59, Consumer Staples 37, Energy 36,
Real Estate 29, Utilities 28, Materials 25, Telecommunications Services 5 (a sector that
was largely reclassified into Communication Services after this membership table was
compiled --- itself a small piece of evidence about the table's vintage).

The cross-sectional median annualised volatility is \textbf{28.95\%} --- high, and
consistent with a large-cap universe measured over a window that includes 2020 and 2022.

\subsection{Survivorship bias: quantified disclosure}
\label{sec:data:survivorship}

The 195 names are those for which Yahoo returns a full or near-full history \emph{today}.
Names delisted, acquired, or bankrupted during 2010--2026 are largely absent, so the panel
conditions on survival to 2026. Consequences, stated plainly:

\begin{enumerate}
  \item \textbf{The universe's historical tradability is overstated.} A trader in 2012
    could not have known which of the 2012 names would still be listed in 2026; our
    backtest implicitly does know.
  \item \textbf{The direction of the bias is upward.} Delistings are disproportionately
    bad outcomes (distress, failed mergers, buyouts at a premium that ends the series).
    Removing them removes both the disasters and the acquisition pops, but for a
    \emph{relative-value} strategy the dominant effect is the removal of permanent
    relationship breaks --- exactly the tail events that destroy pairs trades. Our
    measured drawdowns are therefore optimistic.
  \item \textbf{The membership table is a 2016-vintage snapshot.} Names added to the index
    after 2016 are systematically excluded, which is why the panel is 195 rather than 505
    and why \code{universe\_by\_year} plateaus at $\approx$185 rather than growing.
  \item \textbf{Quantification available to us:} the shortfall itself. The specification
    asked for $\ge200$ names; we deliver 195 with a documented reason
    (\code{qa.py} reports ``195 tickers (spec target 200)''), 96.4\% of which have
    $\ge5$ years of history and 95.9\% $\ge7$ years. Breadth is adequate for the
    four-basket book; it is \emph{not} adequate for a 500-name universe scan, which is
    one reason the multiple-testing audit of Phase~\ref{ch:phase10} covers 44 strategies
    rather than thousands.
\end{enumerate}

\begin{remark}[What a fix would require]
A survivorship-bias-free panel needs (i) point-in-time index membership (yearly constituent
lists, e.g.\ from Wikipedia's historical S\&P 500 change tables or a vendor), joined to
(ii) a price source that retains delisted tickers (CRSP, or a vendor's graveyard file),
and (iii) delisting returns applied on the exit date. None of the three is obtainable in
this sandbox. This is listed first among the non-goals of Chapter~\ref{ch:nongoals}
because it is the limitation most likely to change the conclusions: a survivorship-free
re-run would, we expect, lower every Sharpe reported here and lengthen every drawdown.
\end{remark}

\subsection{Honest reading}

The panel is real, large enough in time (16.7 years, 4{,}190 bars per basket), correctly
adjusted (total return), mechanically clean (10/10 QA), and fully documented with
checksums. It is too narrow in cross-section (195 names, one vintage of membership) and
survivorship-biased. Every later result should be read as \emph{conditional on this
universe}, and the conditional nature is why Phase~\ref{ch:phase10} counts hypotheses
rather than assuming the basket choices were free.

\subsection{Exit gate --- status}

\begin{itemize}
  \item[\checkmark] $\ge200$ names $\times\ge10$ years: \textbf{195} names
    $\times$16.7 years --- name target missed by 5, disclosed and explained;
    year target exceeded by 67\%.
  \item[\checkmark] corporate actions documented and handled (total-return
    \code{adj\_close});
  \item[\checkmark] no lookahead in panel construction (strict balanced export, no
    forward fill);
  \item[\checkmark] \code{results/tables/data\_summary.csv} written by code;
  \item[\checkmark] survivorship analysis in the report.
\end{itemize}
